\subsection{N-017: an actual minimal block, an order-dependent failure, and a unimodality obstruction}
\label{sec:nonstandard-first-block}

\paragraph{Precise scope of the construction.}
This result executes N-016 on an actual rank-four block.  It retains
the additional bar-fixed condition and rational coefficient field of
N-013, the actual unsaturated order $B_S=\mathbb Q[D]\langle P_i\rangle$
of N-015, and the rational GT matrix units of N-010.  It makes two
different, explicit choices for the bitableau order.  Neither choice
is retroactively attributed to VIII.  The ambient coordinates are
the opposite KL tensor coordinates (N-61).

There is an outer-label repair in which $S'_{22}$ is a minimal
label: use inclusion of the dominance lower sets of maximal
ordinary-weight antichains, keeping distinct equal antichains
incomparable, as in N-014.  The ordinary partition chain at size four is
$4>31>22>211>1111$.  The maximal antichain for $S'_{22}$ is
$\{(22,22)\}$, whereas that for the other minimal label
$\theta S'_{31}$ is $\{(31,211),(211,31)\}$.  These antichains are
incomparable.  The remaining maximal antichains, obtained from the
exact tensor decompositions of N-010, contain one of these two
dominance lower sets.  Thus the first block calculation requires no
choices of scaling in any other outer block.

The two paths of $S'_{22}$ are $N_2-$ and $N_2+$, written simply
$-$ and $+$.  We use $-<+$ on paths.  The first choice is the product
order on the row and column paths:
\[
 (-- )<(-+),(+-)<(++),\qquad (-+)\ \hbox{and}\ (+-)\ 
 \hbox{incomparable}.
 \tag{N-78}
\]
The second choice is its row-first total refinement
$--<-+<+-<++$.  Other outer labels can be placed after this block
in a total refinement when that variant is used.

For the ambient degree comparison, a genuine factor order is
specified by $u\preceq w$ if $u=w$ or $\ell(u)>\ell(w)$, and tensor
coordinates have the product order.  Distinct permutations of equal
length are therefore incomparable.  The calculation also passes,
with the same successful elements and the same failed stage, for
the broader convention comparing numerical length pairs and making
equal length pairs incomparable.  This removes any reliance on that
possible interpretation of VIII's reverse-length wording.  The
comparison of degree tuples is exactly VIII's quantified condition
quoted and analyzed in N-016, rather than a single lexicographic test.

\paragraph{The actual module and its verification.}
Put $L=(D-4)^2(D-2)$ and define
\[
 \begin{aligned}
 e&=D L(D-1)^2,\qquad a=D^2e,\qquad b=D^3 L(D-2)(D-1),\\
 c&=D^3(D-4)^2(D-1)^2,\qquad k=(D-2)/D^2.
 \end{aligned}
 \tag{N-79}
\]
In the GT matrix basis ordered $(+,-)$, the intersection of $B_S$
with this full matrix block has the following actual $S$-basis:
\[
 A=\begin{pmatrix}0&0\\0&a\end{pmatrix},\quad
 B=\begin{pmatrix}0&0\\ b/D&2a/D\end{pmatrix},\quad
 C=\begin{pmatrix}0&c/D\\0&a/(2D)\end{pmatrix},\quad
 E=\begin{pmatrix}e&-2c/D^2\\-b/(2D^2)&-(D+1)e\end{pmatrix}.
 \tag{N-80}
\]
The apparent divisions in these displayed formulas are exact
polynomial divisions.  The rows are, respectively, rows
$113,112,111,110$ of the zero-indexed GT Hermite certificate.

Here are the exact checks establishing that (N-80) describes the
unsaturated intersection.  The N-015 order basis was mapped to the
nine GT blocks and reduced by polynomial Euclidean row operations,
in reverse GT order.  All $114$ pivots are present.  A straight-line
program of $11731$ polynomial row operations gives every new row in
the old order basis; exact polynomial triangular division puts each
old row in the new module.  Thus both module containments hold.
The last four rows have GT support exactly in this block.  Conversely,
any element supported there must have zero coefficients on all earlier
Hermite pivots, successively, so these four rows span the intersection.
The same argument gives the smaller path down-set intersections used
below; in particular $\operatorname{span}_S\{C,A\}$ has no hidden
contribution from $B$, since the excluded $-+$ pivot forces its
coefficient to vanish.

The separate verifier
\nolinkurl{verify_order_gt_hermite.py} does not trust the constructed
word-to-GT evaluation.  It verifies that its proposed map sends the
identity to the block identities and intertwines right multiplication
by each $P_i$ on every one of the $114$ actual order basis rows.
These $342$ identities use the standard-$T$ ambient multiplication
reconstructed in N-015's separate verifier.  Since $1$ generates the
word order under right multiplication, these identities determine
the map on every word.  It then replays the whole Hermite program
jointly in GT and ambient coordinates and checks the reverse
containment by polynomial division.  Consequently neither a
generic projection nor a saturation is used in (N-80).

\paragraph{Exact finite spaces in the globalization problem.}
The first three product-order stages have spaces
\[
 \begin{array}{c|c|c|c}
 T&V&\text{primitive current-line generator}&V/(S t)\\\hline
 --&SA&A&0\\
 -+&S(DB-2A)\oplus\mathbb Q B&DB-2A&\mathbb Q B\\
 +-&S(DC-A/2)\oplus\mathbb Q C&DC-A/2&\mathbb Q C.
 \end{array}
 \tag{N-81}
\]
For example, at $-+$ the previously determined exponent of $--$ is
$-8$, so its lower-coordinate degree bound is $7$.  In $pB+rA$ that
coordinate is $a(2p/D+r)$, with $\deg a=8$.  Writing $p=Du+s$ with
$s\in\mathbb Q$, the bound forces $r=-2u$.  This proves the second
row of (N-81) in the actual order; the third row is identical with
$2$ replaced by $1/2$.

For the last stage put
\[
 T_0=D^2E+2DC+(D/2)B+(D-1)A=aE_{++}.
 \tag{N-82}
\]
The lower exponents for $--,-+,+-$ in the product-order variant
are $-8,-8,-7$.  Directly imposing their exact degree bounds gives
\[
 V=S T_0\oplus
        \operatorname{span}_{\mathbb Q}\{E,C,B,DE+A\}.
 \tag{N-83}
\]
For completeness, write an element as $pE+rC+sB+tA$.  After subtracting
a polynomial multiple of $T_0$, its first coefficient can be taken
as $p=x+wD$.  The $+-$ and $-+$ bounds say respectively
$\deg(Dr-2p)\le1$ and $\deg(Ds-p/2)\le1$, forcing $r,s$ constant.
The $--$ bound says
$\deg(D^2t+Dr/2+2Ds-(D+1)p)\le1$, forcing $t=w$.
This is exactly (N-83).  The coefficient $D^2$ in the primitive
current-line vector is minimal: solving the three lower equations
over $F$ gives coefficient ratios with least common denominator
$D^2$, and their cleared polynomial coefficients have gcd one.

In every row of (N-81), and in (N-83), the cutoff $B$ of (N-73) is
zero.  This is checked by comparing the finitely many ambient
polynomial degrees of the explicitly supplied vectors: for each
nonzero coordinate of the primitive current-line vector, every
complement vector has no greater polynomial degree.  It follows
from N-016 that the spaces to test have dimensions $1,2,2,5$,
respectively.  This is an all-degree reduction, not a degree-limited
search.

The exact linear feasibility tests (N-76)--(N-77) give unique least
profiles in those four finite spaces.  The resulting elements and
exponents are
\[
 \begin{array}{c|c|r}
 T&e_T&a_T\\\hline
 --&X=A&-8\\
 -+&Y=DB-2A&-8\\
 +-&Z=DC-A/2&-7\\
 ++&W=E+C+B/4&-6.
 \end{array}
 \tag{N-84}
\]
The complete inputs, all ambient coordinates, current-coordinate
numerators, quotient bases, and exact outcomes are in
\nolinkurl{bar_fixed_block_rank4_verified.json}; the replay is
\nolinkurl{bar_fixed_block_rank4.py}.  It tests the existence of a
source-comparator failure at every nonzero coordinate of the putative
least profile.  Thus the assertion that these profiles are least
does not follow merely from the lexicographic candidate.

The small linear systems have an additional transparent description.
At $-+$ and $+-$, eliminating the coefficient of
$C_{4321}\otimes C_{4321}$ leaves the displayed pure current-line
vector.  At $++$, the only steps increasing the rank of the greedy
constraints occur at coordinate indices $299,221,272$, representing
\[
 (4321,4321),\qquad(2431,4321),\qquad(3421,3421).
\]
On $(E,C,B,A)$ their normalized ambient coefficient rows are
\[
 L(-1,1/2,2,D),\qquad
 L(-1/2,0,2,D/2),\qquad
 (L/2)(1,0,0,0).
 \tag{N-85}
\]
The first two coordinates are set to zero and the last to degree
at most three.  On (N-83) these conditions leave exactly the line
$\mathbb QW$.  All the remaining source-comparison tests pass.

Full bar-fixed uniqueness also holds.  N-013 excludes an odd current
exponent at a minimum.  Any odd homogeneous perturbation is therefore
$\delta w$ with zero current GT coordinate.  At each of the first
three stages the lower-coordinate bounds force $w=0$.  At the last
stage those bounds permit only $w=\alpha B+\beta C$: the $A$ coefficient
would have degree at least eight whereas its bound is seven.
The first two zero coordinates in (N-85) force
$2\alpha+\beta/2=0$ and $2\alpha=0$, hence $w=0$.
The full minimizing space has rational dimension one, so the
leading-coefficient normalization gives precisely one pair $(a_T,e_T)$.

\paragraph{The total refinement has no fourth global minimum.}
For the row-first total order, the first two elements remain $X,Y$.
At its third stage the lower $-+$ coordinate is also allowed, and
the exact finite space is
\[
 \operatorname{span}_{\mathbb Q}\{C,B,DC-A/2\}.
\]
Its unique full-domain minimum is $C-B/4$ with exponent $-6$.
At the fourth stage the different prior exponent changes the exact
space to
\[
 V=S T_0\oplus
          \operatorname{span}_{\mathbb Q}\{E,B,DE+2C+A\}.
 \tag{N-86}
\]
The cutoff remains zero.  The unique lexicographic minimizing line
has normalized representative
\[
 U=(D+2)E+2C+B/2+A,
 \qquad a_U=-7.
 \tag{N-87}
\]
The admissible element $U'=E+B/2$, with exponent $-6$, contradicts
its being a global least profile: at $(3421,3421)$ its coefficient
is $L/2$, whereas that of $U$ is $(D+2)L/2$.  Its degree is strictly
smaller, and at every strict predecessor its degree is at most that
of $U$.  This last finite assertion is checked directly on the full
ambient coordinates in the certificate; the stricter factor-order
verifier also finds a failure already at $(3412,3412)$.
Thus $d(U)\not\le d(U')$ under the actual source comparator.
Every global least profile would have to be lexicographically least,
so no global minimum exists at this stage.  The distinction between
the two choices of order is therefore mathematically consequential.

\paragraph{Multiplication in the successful block.}
The four elements in (N-84) have matrices
\[
 X=aE_{--},\qquad Y=bE_{-+},\qquad Z=cE_{+-},\qquad
 W=\begin{pmatrix}e&kc\\ kb/4&-e\end{pmatrix}.
 \tag{N-88}
\]
They form an $F$-basis of the matrix block.  Their complete $64$
structure constants, of which $28$ are nonzero, are stored in
\nolinkurl{bar_fixed_block_constants_verified.json} and computed
by \nolinkurl{bar_fixed_block_constants.py}.  Each nonzero constant
has the form
\[
 \varepsilon r\,D^i(D-4)^2(D-2)^j(D-1)^l,
 \quad \varepsilon\in\{1,-1\},\quad r\in\mathbb Q_{>0},
 \quad i\ge1,\quad1\le j\le3,\quad l\ge1.
 \tag{N-89}
\]
Here $D-4=(v-v^{-1})^2$, and $D=q+2+q^{-1}$ with $q=v^2$.
For $j=1,2,3$, respectively, the coefficient lists of
$D(D-2)^j$ from their lowest to highest powers of $q$ are
\[
 (1,2,2,2,1),\quad
 (1,2,3,4,3,2,1),\quad
 (1,2,4,6,6,6,4,2,1).
\]
They are symmetric and unimodal.  The same is true of $D$ and
$D-1$.  Products of nonnegative symmetric unimodal Laurent
polynomials are symmetric unimodal, as follows by expressing each
as a nonnegative sum of centered interval polynomials and multiplying
those intervals.  Equation (N-89) therefore proves the claimed
signed positivity and unimodality of every multiplication constant,
after the factor $(v-v^{-1})^4$ is removed.  Coefficients here are
indexed by every integer $q$ exponent, so parity gaps in $v$ are
not silently discarded from a purported statement about consecutive
$v$ exponents.

\paragraph{The ambient coefficients nevertheless fail unimodality.}
The coefficient of $W$ at the diagonal tensor coordinate
$C_{3421}\otimes C_{3421}$ is exactly
\[
 \frac12(D-4)^2(D-2)
   =(v-v^{-1})^4\,\frac{q+q^{-1}}2.
 \tag{N-90}
\]
The residual coefficient sequence is $(1/2,0,1/2)$, which is not
unimodal even when indexed by consecutive integer powers of $q$.
There is no alternative choice of the vanishing factor making it
nonnegative and unimodal.  Indeed a nonzero Laurent polynomial with
nonnegative coefficients has strictly positive value at $v=1$.
The order of vanishing of (N-90) there is exactly four; hence any
factorization $\pm(v-v^{-1})^m f$ with nonnegative Laurent $f$ forces
$m=4$.  Its sign is then positive, and $f=(q+q^{-1})/2$ is forced.

This is a failure of the structural-coefficient unimodality property
for the explicit repaired variant (N-78), even though the
multiplication constants of its four-element block pass the test
in every product.  It does
not identify this variant with a successfully defined printed $E(4)$,
nor assert that every possible repair has this defect.  The full
$114$-element globalization and its full table have not been
computed by this block calculation.

The obstruction is unchanged if the other KL convention is used
consistently.  The semilinear map $v\mapsto v^{-1}$ and
$T_i\mapsto-T_i$ on both tensor factors fixes every $P_i$ and
$\delta$, and sends $C'_w$ to $(-1)^{\ell(w)}C_w$.
It fixes every element of $B_R^{\mathrm{bar}}$ by
$B_R^{\mathrm{bar}}=B_S\oplus\delta B_S$.
Consequently the two ambient coefficient systems differ only by
the corresponding length signs and bar conjugation; their degrees
and the signed unimodality failure agree.

\paragraph{Source and check boundary.}
The original degree comparator and globalization/scaling domain are
in VIII, section 3.3, printed pages 20--22.  The structural-coefficient
positivity claim is in its subsection ``Positivity in the Kronecker
case.''  The exact word-order and GT inputs are N-010 and N-015,
and the finite all-degree justification is N-016.  The accompanying
computations are exact rational polynomial calculations.  The root
agent independently read this entire proof, both stage helper modules,
the construction and constant-table scripts, and the separate GT-map
verifier, then replayed all three exact computations successfully.
The receipt is \nolinkurl{nonstandard_root_review_N017.md}.
